保存ログからの狭い条件付き区間は，Hの欠測２生成に正の生成進捗がなく前提不成立となった．欠測４件は残し，以下は公開レシートと非負性だけを使った事後集計である．
\begin{table}[tb]
\centering\fontsize{8pt}{10pt}\selectfont\setlength{\tabcolsep}{3pt}
\caption{２反復全件のLと各対照との差．下限は欠測計数の非負性だけによる．時間は実測差である．}
\begin{tabular}{lrr}
測定 & L−M & L−H \\ \hline
論理token & $\geq +197,929$ & 未確定 \\
cache除外入力 & $\geq +19,977$ & 未確定 \\
出力token & $\geq +20,204$ & 未確定 \\
作業秒 & +1,702.316 & +2.611 \\
\hline\end{tabular}
\end{table}
L−Mの論理token差は，成功同士5組では$-132,765$，失敗同士2組では$\geq +347,221$，成否不一致1組では$-16,527$であった．各区分は同じ４課題の反復を含み，独立標本ではない．成功部分の減少を全仕事へ拡張せず，Mの成功喪失も保持する．
L−Hは両側に欠測があるため既知部分の差を全費用差へ読み替えない．片側境界は信頼区間ではなく，GPU演算量，消費電力，課金，因果的な高速化も測っていない．
